% =========================================================================
% OVERLEAF RESUME SNIPPET: PROJECTS SECTION
% Copy and paste this directly into your Overleaf LaTeX editor under PROJECTS
% =========================================================================

\section*{PROJECTS}
\vspace{-4pt}
\hrule
\vspace{6pt}

\begin{itemize}[leftmargin=12pt, label=\textbullet]

    % PROJECT 1: FlexCharge (Flagship Project)
    \item \textbf{FlexCharge: Level-3 DC Fast Charging Station with V2G Grid Optimization} \hfill \textit{Oct. 2026} \\
    \textit{Modeled a 50 kW bidirectional EV charging station with active grid support and 800V DC bus stabilization.} \hfill \href{https://github.com/rakshitsh45/FlexCharge}{\underline{Github}}
    \begin{itemize}[leftmargin=15pt, label=-]
        \item \textbf{Tools \& technologies used}: MATLAB, Simulink, Simscape Electrical, Synchronous $d$-$q$ Vector Control, SRF-PLL, State-Space Modeling, IEEE 519-2022, IEEE 1547-2018.
        \item \textbf{Architected a 3-stage power structure} comprising a 3-phase Active Front End (AFE) PWM converter, a regulated 800V common DC bus ($C_{dc} = 4700\ \mu\text{F}$), and a bidirectional Buck-Boost DC-DC stage interfacing a 60 kWh Li-Ion pack.
        \item \textbf{Designed a dual-loop decoupled $d$-$q$ vector controller} with Symmetrical \& Modulus Optimum PI tuning, achieving \textbf{Unity Power Factor ($\cos\phi = 1.000$)} in G2V mode and holding DC bus deviations below $\pm 0.18\%$ under 50 kW step reversals.
        \item \textbf{Engineered dynamic voltage sag support (IEEE 1547)} injecting \textbf{$+17.89$ kVAR reactive power} within 1.8 ms of a 12\% grid sag, and implemented \textbf{V2G peak shaving} delivering 22.2 kW to the grid with \textbf{1.71\% THD} (IEEE 519 compliant).
    \end{itemize}

\end{itemize}
